% ==================================================================
%  CLUSTERING METHOD TEMPLATE --- copy me, do not edit in place.
%
%    cp template/method_template.tex \
%       documentation/sections/clustering_methods/<family>/NN-<name>.tex
%
%  Then: replace <NAME>, replace the label suffix `template' with the
%  method's short name, and add one \input line for the new file under
%  the matching family subsection in cluster_main.tex.
%
%  This file is NOT included in the build.
%
%  Levels: a method is a \subsubsection sitting under its family
%  subsection in cluster_overview.tex, so the parts below are
%  \paragraph --- the article class's fourth level (there is no
%  \subsubsubsection).  Keep their order unchanged in every copy: the
%  comparison section (Section~\ref{sec:comparison}) relies on every
%  method reporting the same things in the same order.
%
% ------------------------------------------------------------------
%  REFERENCES --- a source is declared, not merely cited.
%
%  1. Add every source you use to the shared citation key mapping
%     sheet.  Everyone declares their references there and nowhere
%     else:
%     https://deakin365-my.sharepoint.com/:x:/g/personal/s224236373_deakin_edu_au/IQDSW3PBbHhZSq6vzGGQqi3WAbgJ1id1-t02k5B6fduX5Ys?e=nShY5e
%  2. Each entry in the sheet carries a key of the form ref_<n>.  Cite
%     it in the text as \cite{ref_<n>}.  Do not invent your own key,
%     and do not edit documentation/literature.bib yourself --- the
%     document lead applies the mapping there.
%  3. Record every source in the REFERENCES USED block at the end of
%     this file, including anything you read but did not cite.
%
%  Cite in the part where the claim is made, not in a lump at the end.
%  Every formulation, every published result and every parameter value
%  taken from the literature carries its own citation.
%
%  ACADEMIC INTEGRITY --- an uncited claim, figure, equation, table or
%  code fragment is plagiarism, whether it came from a paper, a
%  library's documentation, a teammate, or a generative AI tool.
%  Quote and cite it, or write it yourself.  Declare AI assistance as
%  Deakin policy requires.
%
%  Before opening the pull request, run your section through Turnitin.
%  Do not commit the report --- the repository is not where it belongs.
%  Paste the similarity summary into the pull request as an image, or
%  send it to the document lead privately; the lead may also run the
%  check later.  Policy and procedure:
%     https://d2l.deakin.edu.au/d2l/le/content/93067/viewContent/5882569/View
% ==================================================================
\subsubsection{<NAME>}
\label{sec:tech:template}

\paragraph{Overview}
\label{sec:tech:template:overview}
\placeholder{What the algorithm does, in one paragraph, and the family
it belongs to (see Section~\ref{sec:background:families}).}

\paragraph{Assumptions and applicability}
\label{sec:tech:template:assumptions}
\placeholder{What the algorithm assumes about cluster shape, density,
scale and noise; the conditions under which it is a sensible choice
for the data described in Section~\ref{sec:data}.}

\paragraph{Formulation}
\label{sec:tech:template:formulation}
\placeholder{Objective function and update rule, using the symbols
declared in the notation table. Cite the original source.}

\paragraph{Hyperparameters and tuning}
\label{sec:tech:template:hyperparameters}
\placeholder{Each hyperparameter, the range searched, and the value
selected --- with the reason.}

\paragraph{Complexity}
\label{sec:tech:template:complexity}
\placeholder{Time and space complexity, and what that means at
\projectname{} data volumes.}

\paragraph{Application to \projectname{} data}
\label{sec:tech:template:application}
\placeholder{Implementation, any deviations from the shared
preprocessing pipeline, and why they were necessary.}

\paragraph{Results}
\label{sec:tech:template:results}
\placeholder{Metrics from Section~\ref{sec:methodology:metrics},
cluster profiles, and supporting figures.}

\paragraph{Limitations}
\label{sec:tech:template:limitations}
\placeholder{Where this method failed or would be expected to fail
on this data.}

\paragraph{Summary}
\label{sec:tech:template:summary}
\placeholder{Two or three sentences: the verdict on this method,
carried into the comparison section.}

% ==================================================================
%  REFERENCES USED --- fill this in.  It is checked at review.
%
%  One line per source, using the ref_<n> key from the citation key
%  mapping sheet.  Include what you read and did not cite: a reviewer
%  needs to see what the write-up rests on, not only what it quotes.
%
%    ref_12   Author (Year), short title       cited in: Formulation
%    ref_13   Author (Year), short title       read, not cited
%
%  These same keys go in the `references' field when the matching
%  component is registered in xxcluster --- see CONTRIBUTING.md, 2.3.1.
%  The two lists must agree.
% ==================================================================
